\section{Molecular dynamics}
\subsection{Basic statistical mechanics}
Consider a box consider N-partciles, we would like to study the property of this system. The particles in this system has interaction. In MD, we have a force field and the field here depends on position. In this course, we mainly only consider the Coulomb's interaction:
\begin{equation}
    F=k\frac{q_1 q_2}{r^2}
\end{equation}

The approach in classical mechanics includes:\\
1. Newton's equation of motion,

\begin{equation}
    \textbf{F}=m\frac{d^2 \textbf{x}}{dt^2}
\end{equation}

2: Hamiltonian mechanics.
\begin{equation}
    \dot{\hat{p}}=-\frac{\partial H}{\partial q}
\end{equation}
\begin{equation}
    \dot{\hat{q}}=\frac{\partial H}{\partial p}
\end{equation}

3: Lagrangian mechanics

If we want to describe a state, we can have a state point, where the state function $S=(q,p)=(r_1,r_2,...,r_N,p_1,p_2,...,p_N)$. The phase space of the system has $6N$ dimension.

\subsubsection{NVE ensemble}
In NVE ensemble system, the number of particle(N), volume(V), and Energy(E) is unchange in the system. The expected value of A (such as pressure) can be written as:
\begin{equation}
    <A>=\frac{\sum_S\delta_{\varepsilon(S),E} A(S)}{\sum_S \delta_{\varepsilon(S),E}}
\end{equation}
\subsubsection{Canonical ensemble/NVT}
If the Temperature of the system is confined instead of energy, we will have the NVT ensemble.
\begin{equation}
    <A>=\frac{\sum_S e^{-\beta \varepsilon(S)} A(S)}{\sum_S e^{-\beta \varepsilon(S)}}
\end{equation}

Classically, we can define temperature as
\begin{equation}
    \sum^N_{i=1} \frac{\textbf{p}_i^2}{2m} = \frac{3Nk_BT}{2}
\end{equation}
\subsection{Ergodicity}
The Ergodic hypothesis means that the ensemble average can be obtained by the time average:

\begin{equation}\label{eqn:Ergodicity}
    <\textbf{A}>=\lim_{t\rightarrow\infty}\frac{1}{t}\int_0^\infty dt' \textbf{A}(t')
\end{equation}

In MD, the hypothesis means that if the simulation, equation \eqref{eqn:Ergodicity} will give us an accurate results. This works for fluids, supercritical fluids, or gas, but does not work for solid and glass.

\subsection{Molecular Dynamics (MD)}
To perform MD, we need 1. State function $S_i$, and 2. $F_i$: There are two types of Force field $F_i$, including classical MD, and quantum mechanics MD (first-principles MD, or ab initio MD). 3. Given a time steps h, we solve the equation motion of $t\rightarrow t+h$. Step 2 and Step 3 forms a loop and repeat. Step 3 is performed by Verlet method.

\subsection{Verlet method}
Taylor expansion gives us:
\begin{equation}\label{eqn:verlet1}
    x(t+h) = x(t) + \dot{x}(t)h+\frac{1}{2!}\ddot{x}(t)+\frac{1}{3!}\dddot{x}(t)+O(h^4)
\end{equation}
\begin{equation}\label{eqn:verlet2}
    x(t-h) = x(t) - \dot{x}(t)h+\frac{1}{2!}\ddot{x}(t)-\frac{1}{3!}\dddot{x}(t)+O(h^4)
\end{equation}
With some algebra, we have
\begin{equation}\label{eqn:verlet3}
    x(t+h)-x(t-h)=2x(t)+\ddot{x}(t)h^2+O(h^4)
\end{equation}
\begin{equation}\label{equation:Verlet}
    x(t+h)=2x(t)+x(t-h)+\frac{f(t)}{m}h^2+O(h^4)
\end{equation}
Equation \eqref{equation:Verlet} gives us the Verlet equation.

\subsubsection{Time reversal and energy conservation}
Why can't we use the numerical method to solve the PDE? The fundamental reason is due to time reversal symmetry. By taking $t\rightarrow -t$, only $v$ has the sign being reverse, and \eqref{equation:Verlet} remains the same. Therefore, using Verlet method allows us to satisft time-reversal symmetry, and \textbf{allowing energy to be conserved}.

\subsubsection{Velocity in the Verlet method}
By subtracting equation \eqref{eqn:verlet1} by equation \eqref{eqn:verlet2}, we get:
\begin{equation}\label{eqn:vverlet}
    v(t)=\dot{x}(t)=\frac{x(t+h)-x(t-h)}{2h}+O(h^2)
\end{equation}

To perform Velocity Verlet method, we substitue \eqref{eqn:vverlet} to \eqref{equation:Verlet} to get:
\begin{equation}
    x(t+h)=x(t)+hv(t)+\frac{f(t)}{2m}h^2
\end{equation}
This equation allows us to update $x(t+h)$

To update $v(t)$, we replace t by t+h in \eqref{equation:Verlet} to get:
\begin{align}
    v(t+h)&=\frac{x(t+2h)-x(t)}{2h}\\
    &=\frac{x(t+h)+hx(t+h)+\frac{f(t+h)}{2m}h^2}{2h}\\
    v(t+h)&=v(t)+\frac{f(t)+f(t+h)}{2m}h \label{eqn:velocityupdate}
\end{align}

It is easy to see that equation \eqref{eqn:velocityupdate} satisfy time reversal symmetry and \textbf{the proof is left as an exercise for the reader}.

\subsection{Leapfrog method}
Using Leapfrog method can get the exact same result as the Verlet. Leapfrog method is a method similar to R-K method,  but satisfy time-reversal symmetry.

\subsubsection{Apply the leapfrog method in Newton’s equations of motion}
Details are in the lecture notes.

\subsubsection{Can we trust the trajectory obtained by the leapfrog/Verlet method?}
Yes. Even though there are chaotic conditions, Verlet method forced to minimize the action.

How can we compute functional derivative in computer? Consider an action:
\begin{align}
    S=\int^b_a \bigg[\frac{1}{2}m\bigg(\frac{dx}{dt}\bigg)^2-V(x)\bigg] dt\\
    S=\sum^{N-1}_{i=0}\bigg[\frac{1}{2}m\bigg(\frac{x_{i+1}-x_{i}}{h}\bigg)^2-V(x_i)\bigg]h
\end{align}

For $\forall i = 0,1,2,...,N-1$, we have
\begin{align}
    \frac{\partial S}{\partial x_i}=0\\
    \frac{\delta S}{\delta x(t)}=0\\
    x_{i+1}=2x_i=x_{i-1}+f_i\frac{h^2}{m}
\end{align}

\newpage
